\section{多种知识点结合的问题}
\subsection*{话题 3-A：微分与不动点}
\thm{7.3.1（压缩映射原理）}设 $(X,\rho)$ 为完备度量空间，$T:X\to X$ 满足 $\rho(Tx,Ty)\le k\rho(x,y)$，其中 $0<k<1$。则 $T$ 有唯一不动点。
\pf 从任意 $x_0\in X$ 出发，令 $x_{n+1}=Tx_n$。则 $\rho(x_{n+1},x_n)\le k^n\rho(x_1,x_0)$，进而
\[\rho(x_m,x_n)\le\frac{k^n}{1-k}\rho(x_1,x_0)\quad(m>n).\]
完备性给出极限 $x_*$；由 $T$ 的连续性，$Tx_*=x_*$。若 $Ty_*=y_*$，则 $\rho(x_*,y_*)\le k\rho(x_*,y_*)$，所以两者相同。
\thm{7.3.2}设 $\Omega\subset\R^n$ 为凸区域或凸闭区域，$F:\overline\Omega\to\R^m$ 连续并在 $\Omega$ 内可微。则 $F$ 满足
\[|F(x)-F(y)|\le M|x-y|\quad(x,y\in\overline\Omega)\]
的充要条件是 $\|F_x(x)\|\le M$ 对所有 $x\in\Omega$ 成立，其中 $F_x$ 为 Jacobi 矩阵，范数为算子范数。
\pf 对单位向量 $\nu\in\R^m$，对标量函数 $\nu\cdot F$ 沿连接 $x,y$ 的线段使用中值定理，得到充分性；必要性由差商极限得到。
\par 若直接对向量函数套用一维微分中值定理，可能误以为存在同一个 $\xi$ 使 $F(y)-F(x)=F_x(\xi)(y-x)$；一般不存在这样的 $\xi$。先与任意单位向量 $\nu$ 取内积，再对标量函数应用中值定理，就避开了这一问题。最后利用
\[|F(y)-F(x)|=\sup_{|\nu|=1}\nu\cdot(F(y)-F(x))\]
恢复向量范数估计。矩阵范数 $\|F_x\|$ 是对所有单位方向的导数变化取最大值。
\par 因而若 $F$ 把凸闭集映到自身，且 $\sup\|F_x\|<1$，则 $F$ 为压缩映射，其迭代收敛到唯一不动点。

\ex{7.3.2}若 $E\subset\R^n$ 为凸区域或凸闭区域，$F:E\to E$ 可微，且 $\|F_x\|<1$，则对不同的 $x,y\in E$ 有 $|F(x)-F(y)|<|x-y|$。
\pf 对固定线段 $[x,y]$，沿线段的导数范数连续且最大值严格小于 $1$，再使用定理 7.3.2。
\ex{7.3.3}在例 7.3.2 的条件下，对任意内部点 $x_0$ 和以它为中心的紧子集 $W\subset E$，存在 $k<1$，使
\[|F(x)-F(x_0)|\le k|x-x_0|\quad(x\in W).\]
\pf 在 $x_0$ 邻域中对 $\|F_x\|$ 取上界。对于一维区间，也可把 $(F(x)-F(x_0))/(x-x_0)$ 在 $x_0$ 处连续延拓，再取最大值。
\ex{7.3.4}设 $f:\R\to\R$ 可微、有界，且 $|f'|<1$。则 $f$ 有唯一不动点 $\bar x$；从任意 $x_0$ 出发的 $x_{n+1}=f(x_n)$ 均收敛于 $\bar x$。
\pf $f(x)-x$ 在两端分别趋向正、负无穷，故存在零点；例 7.3.2 保证唯一性。迭代轨道有界，例 7.3.3 在轨道闭包上给出统一 $k<1$，从而收敛。
\par 此处 $|f'(x)|<1$ 是逐点严格不等式，不能未经说明就写出 $\sup_{x\in\R}|f'(x)|<1$。轨道有界使我们只需在一个紧区间上取统一常数。一般地，利用“局部严格压缩+迭代有界”替代“全局统一压缩”，正是例 7.3.3--7.3.6 的共同思路。
\ex{7.3.5}设 $E\subset\R^n$ 为凸有界闭区域，$F:E\to E$ 可微，$\|F_x\|<1$，且 $F$ 无不动点。证明存在 $k<1$ 使
\[|F^2(x)-F(x)|\le k|F(x)-x|,\qquad x\in E.\]
\pf 比值函数在紧集 $E$ 上连续，由例 7.3.2 严格小于 $1$，故最大值也小于 $1$。
\ex{7.3.6}设 $F:\R^n\to\R^n$ 可微、$\|F_x\|<1$。若某个 $x_0$ 的迭代轨道 $\{F^k(x_0)\}$ 有界，则它收敛于 $F$ 的唯一不动点。
\pf 若无不动点，在轨道闭包上应用例 7.3.5 的比值估计，轨道成为 Cauchy 列而收敛于不动点，矛盾。存在后再用例 7.3.3 证明收敛。
\par 具体地，若 $F$ 无不动点，紧轨道闭包上有 $k<1$ 使
\[|F^{j+1}(x_0)-F^j(x_0)|\le k|F^j(x_0)-F^{j-1}(x_0)|.\]
逐次迭代得相邻点距离受几何级数控制，从而轨道为 Cauchy 列。其极限 $\xi$ 满足 $F(\xi)=\xi$，矛盾。这个证明只使用一个轨道的有界性，没有预设整个映射空间有界。

\subsection*{话题 3-B：微分与一致收敛}
\ex{7.3.7}设 $f_n$ 在 $[a,b]$ 上可微，导函数一致有界，且 $\{f_n(a)\}$ 有界。则 $\{f_n\}$ 有子列在 $[a,b]$ 上一致收敛。
\pf 导数的一致有界性给出等度连续；加上锚点处有界，推出一致有界。应用 Arzelà--Ascoli 定理。
\ex{7.3.8}沿用例 7.3.7 的条件。证明：(i) $f_n$ 一致有界；(ii) 满足统一 Lipschitz 条件；(iii) 可选子列在 $[a,b]\cap\mathbb Q$ 上收敛；(iv) 此子列在整个 $[a,b]$ 上收敛；(v) 收敛是一致的。
\pf 前两步由中值定理；第三步对可数稠密集作对角选取；第四、五步由一致 Lipschitz 条件及有限网格逼近。
\par 为说明最后一步，给定 $\varepsilon>0$，在 $[a,b]$ 中选有限个有理点 $q_1,\dots,q_s$，使每个 $x$ 距某个 $q_j$ 不超过 $\varepsilon/(3M)$。所选子列在这有限个点上同时成为 Cauchy 列。若 $m,n$ 足够大，
\[|f_m(x)-f_n(x)|\le |f_m(x)-f_m(q_j)|+|f_m(q_j)-f_n(q_j)|+|f_n(q_j)-f_n(x)|<\varepsilon.\]
这就把稠密集上的逐点收敛提升为整个区间上的一致收敛。
\ex{7.3.9}设 $\Omega\subset\R^n$ 为区域、$E\Subset\Omega$，$F:\Omega\to\R^m$ 可微。若 $F_x$ 在 $\Omega$ 上有界，则存在 $M$，对所有 $x,y\in E$ 有
\[|F(x)-F(y)|\le M|x-y|.\]
\pf 在 $E$ 周围取有限个小凸邻域。近距离点可用定理 7.3.2；相隔一定距离的点由 $F$ 在 $E$ 上有界控制。
\ex{7.3.10}设 $f$ 在 $(a,b)$ 处处可导。证明：\emph{存在} $(a,b)$ 的一个子区间 $[\alpha,\beta]$，使 $f'$ 在该子区间上有界。
\pf 反设每个子区间上的 $f'$ 都无界。依次构造闭区间套 $I_{n+1}\Subset I_n$，使 $|I_{n+1}|\le|I_n|/2$，并使相应端点割线斜率绝对值大于 $n$。公共点 $\xi$ 处的差商就不可能收敛到有限的 $f'(\xi)$，矛盾。这里的量词是“存在一个子区间”，不能误读为“任一子区间”。
\par 实际构造时，应让高斜率割线的两个端点逐步逼近公共点。因为若 $f'$ 在某个开区间内无界，中值定理保证该区间中存在绝对值任意大的割线斜率。将这些割线包含在逐层收缩的闭区间中，就得到与 $\xi$ 处可导性冲突的差商列。

\subsection*{话题 3-C：隐函数、参数方程与迭代}
\ex{7.3.14}求 $x+\ln(1+x)+y+\arctan y=0$ 确定的隐函数 $y$ 的 Maclaurin 级数。
\sol 设 $h(t)=t+\ln(1+t)$，$q(t)=t+\arctan t$，则 $q(y)=-h(x)$，$q'(0)=2\ne0$，故 $y=q^{-1}(-h(x))$ 在原点附近解析。由 Lagrange 反演公式，可将\emph{整条} Maclaurin 级数写成
\[y(x)=\sum_{m=1}^{\infty}\frac{[-h(x)]^m}{m}[t^{m-1}]\left(\frac{t}{t+\arctan t}\right)^m.\]
其中 $[t^{m-1}]$ 表示取幂级数中 $t^{m-1}$ 的系数。逐项代入可求任意阶；前三项为
\[y=-x+\frac14x^2-\frac13x^3+O(x^4).\]
\ex{7.3.15}求 $e^y+y=x(1-x)$ 确定的隐函数 $y$ 在 $[0,1]$ 上的最大、最小值。
\sol 左侧对 $y$ 严格递增。设 $a$ 为 $e^a+a=0$ 的唯一解，$b$ 为 $e^b+b=1/4$ 的唯一解。因为 $x(1-x)$ 在端点最小、在 $x=1/2$ 最大，故 $\min y=a$、$\max y=b$。
\ex{7.3.16}求 $e^y+xy=x^2+x+4$ 确定的隐函数 $y$ 在 $[0,1]$ 上的最大、最小值。
\sol 隐函数唯一且 $(e^y+x)y'+y=2x+1$。设 $a$ 为 $e^a+a=6$ 的唯一解，则 $y(0)=\ln4$、$y(1)=a$。驻点 $\xi$ 满足 $y(\xi)=2\xi+1$，代入原方程得 $e^{2\xi+1}=4-\xi^2$；该式在 $(0,1)$ 有唯一解。令 $b=2\xi+1$，则 $e^b+(b-1)^2/4=4$。比较端点与驻点得最大值 $a$，最小值 $b$。
\ex{7.3.17}设 $x=e^t+t$、$y=\cos t-3t$ 确定 $y=y(x)$。(1) 求 $y''$；(2) 求 $\int y(x)\dd x$；(3) 固定 $x_0\ne1$，令 $x_{n+1}=y(x_n)$，讨论有界性与收敛性。
\sol 因 $x'(t)=e^t+1>0$，
\[y'=-\frac{\sin t+3}{e^t+1},\quad y''=\frac{e^t\sin t-e^t\cos t+3e^t-\cos t}{(e^t+1)^3}.\]
又由换元，
\[\int y(x)\dd x=\frac{e^t}{2}(\cos t+\sin t)-3(t-1)e^t+\sin t-\frac32t^2+C.\]
$y$ 严格递减，$y(1)=1$，而 $y'(1)=-3/2$。若轨道收敛，极限只能是 $1$；但在 $1$ 邻域内 $|y(x)-1|\ge(5/4)|x-1|$，矛盾。另一方面，$y(x)/x\to0$（$x\to+\infty$），$y(x)/x\to-3$（$x\to-\infty$），故二次迭代在足够大的正半轴把 $x$ 拉回，轨道有界。

\subsection*{话题 3-D：题目的再加工}
\ex{7.3.18}讨论级数 $\sum_{n=1}^\infty\sin(2n!\pi e)$ 的敛散性。
\sol 令 $N_n=n!\sum_{k=0}^n1/k!\in\mathbb Z$，$r_n=n!e-N_n=1/(n+1)+O(n^{-2})$。则 $\sin(2n!\pi e)=\sin(2\pi r_n)\sim2\pi/n>0$，故级数发散。
\ex{7.3.19}讨论 $\sum_{n=1}^\infty\sin(n!\pi e)$ 的敛散性。
\sol $N_n\equiv n+1\pmod2$（$n\ge2$），故 $\sin(n!\pi e)=(-1)^{n+1}\sin(\pi r_n)$。$r_n\downarrow0$ 且 $\sin(\pi r_n)\sim\pi/n$，因此级数条件收敛。\textbf{校勘：}原书把例 7.3.18 写成交错收敛，但按所印的 $2n!\pi e$，其项最终为正且与 $2\pi/n$ 等价；两题的结论应如上。
\ex{7.3.20}设 $E$ 是级数 $\sum_{n\ge1}\sin(n!\pi x)$ 的收敛点集。讨论 $E$ 在 $\R$ 中是否稠密、是否可含无穷多个无理数、其余集是否可含无穷多个无理数，并找出尽可能多的具体点。
\sol 若 $x$ 为有理数，则当 $n$ 足够大时 $n!x$ 为整数，故所有有理数属于 $E$，$E$ 稠密。无理点方面，例 7.3.18--7.3.19 立即给出：$e\in E$，而 $2e\notin E$。更一般地，奇数倍 $e$ 对应交错的 $1/n$ 量级项，属于 $E$；非零偶数倍 $e$ 对应同号的 $1/n$ 量级项，不属于 $E$。故两侧都有无穷多个无理数。原书“对任何整数 $k$，$2ke$ 处发散”应排除 $k=0$：在 $x=0$ 时每项均为零。
\par 关于 $E$ 的余集，单举例子不足以回答“是否有无穷多个无理数”的一般性问题；上述 $2ke$（$k\ne0$）已经给出一族。任意实数还可借助阶乘进位展开控制 $n!x$ 的小数部分，但这部分属于开放式探究，原书也未给完整分类。
\ex{7.3.21}比较 $\int_0^1\!1/x\dd x$、$\iint_{[0,1]^2}\!1/(x+y)\dd x\dd y$、$\iiint_{[0,1]^3}\!1/(x+y+z)\dd x\dd y\dd z$ 的敛散性。
\sol 第一式发散；后两式收敛。特别地，二维积分按 $y$ 先积得到 $\int_0^1\ln((x+1)/x)\dd x=2\ln2$。三维情形在原点附近按径向估计 $r^{-1}$，体积元为 $O(r^2)\dd r$，故可积。
\ex{7.3.22}讨论下列广义积分的敛散性；收敛时求值：
\begin{enumerate}[label=(\roman*)]
\item $\displaystyle\int_0^1\frac{\dd x}{\ln x}$；
\item $\displaystyle\iint_{[0,1]^2}\frac{\dd x\dd y}{\ln x+\ln y}$；
\item $\displaystyle\iiint_{[0,1]^3}\frac{\dd x\dd y\dd z}{\ln x+\ln y+\ln z}$；
\item $\displaystyle\iiint_{[0,1]^3}\frac{\dd x\dd y\dd z}{\ln x+2\ln y+3\ln z}$；
\item $\displaystyle\iint_{[0,1]^2}\frac{\dd x\dd y}{p\ln x+q\ln y}$，$p,q>0$。
\end{enumerate}
\sol (i) 在 $x=1$ 附近发散至 $-\infty$。(ii)--(v) 令 $x=e^{-u}$、$y=e^{-v}$、$z=e^{-w}$，并用 $1/s=\int_0^\infty e^{-st}\dd t$，分别化为
\[-\int_0^\infty\frac{\dd t}{(1+t)^2}=-1,\qquad
-\int_0^\infty\frac{\dd t}{(1+t)^3}=-\frac12,\]
\[-\int_0^\infty\frac{\dd t}{(1+t)(1+2t)(1+3t)}=2\ln2-\frac32\ln3,\]
\[-\int_0^\infty\frac{\dd t}{(1+pt)(1+qt)}=
\begin{cases}-\dfrac{\ln(q/p)}{q-p},&p\ne q,\\-\dfrac1p,&p=q.\end{cases}\]
